\documentclass[conference]{IEEEtran}
\IEEEoverridecommandlockouts
\usepackage{cite}
\usepackage{amsmath,amssymb,amsfonts}
\usepackage{graphicx}
\usepackage{textcomp}
\usepackage{xcolor}
\usepackage{booktabs}
\usepackage{multirow}
\usepackage{array}
\usepackage{url}

% ---- convenience macros -------------------------------------------------
\newcommand{\todo}[1]{\textcolor{red}{\textbf{[TODO(DATA): #1]}}}
\newcommand{\fbetaoof}{\texttt{fbeta\_oof}}

\begin{document}

\title{A Physical S7comm Intrusion Detection Dataset\\
with Session-Shift Evaluation and a Process-Logging Audit}

\author{
\IEEEauthorblockN{1\textsuperscript{st} Thai Thi Minh Tu, 2\textsuperscript{nd} Vuong Linh Thao, 3\textsuperscript{rd} Dang Le Dinh Trang}
\IEEEauthorblockA{
Le Quy Don Technical University, Hanoi, Vietnam\\
thaithiminhtu.338@gmail.com, thaolinh252512@gmail.com, trangdld@lqdtu.edu.vn}
}
% NOTE: remove/blank the author block above if the venue requires double-blind review.

\maketitle

\begin{abstract}
Intrusion-detection research for Siemens S7comm is constrained by a shortage of datasets collected on
physical hardware that pair network observations with process state and report evaluation under a
distribution shift rather than in-distribution cross-validation alone. We present a physical S7comm
dataset built on an S7-1500 testbed running two industrial control applications, aligning network
capture, PLC tag/process logs, and an attack-event timeline by timestamp into 55{,}902 labeled 2-second
windows (47{,}460 benign, 8{,}442 attack) across a six-day collection schedule, organized through a
four-layer feature taxonomy: network volume, ICS protocol presence, S7 operation semantics, and process
state. We evaluate four baselines under three protocols: episode-grouped cross-validation, a Day-5
temporal holdout, and a randomized-order, reduced-rate Day-6 session, treated as exploratory rather than
confirmatory because its anomalous early result triggered a data-quality audit of our own pipeline. Under
a grouped out-of-fold threshold calibration, network-only binary Macro-F1 falls from 0.83--0.91
in-distribution to 0.60--0.66 on Day 6, and the drop is model-dependent. Naive network/process fusion
shows no consistent binary effect across models. We trace this to a concrete collection defect: the
process logger was started manually and overlapped only part of each attack period, so only 839 attack
windows (9.9\%) carry a discriminative process signal, and a matched comparison restricted to real
process observations is structurally uninformative on this release rather than merely unrun. We report
this diagnosis in full, together with the leakage it induced in an earlier pipeline
version, because it determines what the fusion result can and cannot support.
\end{abstract}

\begin{IEEEkeywords}
industrial control systems, S7comm, intrusion detection dataset, PLC testbed, distribution shift, data quality
\end{IEEEkeywords}

\section{Introduction}

Industrial Control Systems built around Siemens S7-family PLCs remain widely deployed in manufacturing,
utilities, and distributed infrastructure. The S7comm protocol governing HMI--PLC and engineering-station
communication was not designed with authentication or encryption by default. A 2026 taxonomic review of 23
public ICS/IIoT intrusion-detection datasets finds protocol coverage heavily skewed toward
Modbus---appearing in 13 of 23 datasets against S7comm in only five---and that most datasets model short,
isolated attack episodes rather than a reconnaissance-through-impact progression \cite{b1}.

Three gaps follow. First, S7comm is under-served relative to its deployment footprint. Second, datasets
that do target it are typically collected on a single control process, making it impossible to tell
whether a detector has learned protocol behavior or one program's idiosyncrasies. Third, published
results are usually reported under in-distribution cross-validation, which systematically overstates what
a detector will do on traffic it has not seen.

The third gap motivates this paper. We designed the collection schedule around an explicit
session-disjoint split from the outset and evaluate every model under every protocol rather than
reporting the best model per split. Days 1--6 share the same PLC, attacker host, and network topology, so
the shift we measure is a within-testbed, session/schedule-level shift (day, episode order, attack
rate)---not cross-testbed or cross-deployment generalization. We use this narrower framing throughout
rather than the broader term ``out-of-distribution.''

Existing S7comm work has focused on protocol modeling for automaton-based detection \cite{b2}, a physical
mining/refinery testbed with a fixed process \cite{b3}, or synthetic attacks injected into a
single-process capture \cite{b4}. What remains rare is a dataset that preserves the relationship between
the sender, the command, the targeted memory area, and the resulting process-state change---since in
S7comm a single correctly addressed write can corrupt a control cycle without any protocol-level
anomaly---while exercising more than one control application and reporting session-disjoint evaluation.

We organize the work around three research questions:
\begin{itemize}
\item \textbf{RQ1 (Coverage).} What operating modes and attack behaviors are represented across the
six-day collection schedule of a physical S7comm testbed, and how are they distributed by day and class?
\item \textbf{RQ2 (Semantics and data quality).} What detection signal is contributed by each feature
layer, and how is that assessment affected by incomplete process-log coverage?
\item \textbf{RQ3 (Generalization).} How does performance change from episode-grouped in-distribution
evaluation to temporally disjoint, randomized-order evaluation under a within-testbed schedule shift?
\end{itemize}

Our contributions are: (i) a physical S7-1500 testbed with ERSPAN mirroring, exercising two instrumented
control applications; (ii) three timestamp-aligned data sources merged into 55{,}902 labeled windows with
explicit leakage controls; (iii) a symmetric benchmark of four baselines across three evaluation
protocols, including a layer-wise ablation; and (iv) an audited finding that network/process fusion
provides no consistent benefit, in which we identify the collection defect responsible rather than
reporting the fused score at face value. We keep the multi-application claim narrow: by collection design
both applications contribute windows to every day of the schedule, though the current release has no
application-id field to audit this per window (Section~\ref{sec:limitations}), and we do not report a
leave-one-application-out evaluation; this establishes only that the dataset is not tied to a single
process, not that a detector transfers across control programs. We position this as a dataset and
benchmark contribution, not a state-of-the-art detector.

\section{Related Work}

\textbf{S7comm protocol modeling.} Kleinmann and Wool model the periodic nature of S7 traffic as a
per-channel Deterministic Finite Automaton, flagging out-of-sequence messages and unexpected memory
references \cite{b2}. This establishes structural regularity exploitable for anomaly detection, but does
not couple network observations to process-state change, nor release a labeled, scenario-organized
dataset.

\textbf{Physical S7comm testbeds.} Rodofile et al.\ release QUT\_S7comm from a physical mining/refinery
testbed, with roughly nine hours of traffic and process logs across 13 attack types \cite{b3}---the
closest prior work in spirit, but targeting a single process without session-disjoint evaluation.
Al-Duwairi et al.\ release SCADA captures from an S7-1500-based medical waste incinerator with synthetic
attacks injected post hoc rather than executed live \cite{b4}. Kellerer et al.\ release a physical-testbed
S7comm dataset spanning two process contexts---a battery storage plant on an S7-1512 and a PV plant on an
S7-1516---with live protocol-level write attacks \cite{b10,b11}.

\textbf{Multi-protocol and multimodal ICS datasets.} G\'omez et al.\ build the Electra dataset from an
electric traction substation using Modbus TCP and S7comm \cite{b5}. Dehlaghi-Ghadim et al.\ present an
anomaly-detection dataset on a heterogeneous-protocol physical testbed \cite{b6}. Shafir et al.\ present
ICS-MMOC3, a multimodal S7comm dataset from a power-generation ICS pairing operational faults with
cyber-attacks \cite{b12}. Zhou et al.\ present ICS-NAD, spanning three ICS vendor platforms and two
industrial processes with 20 attack types and released PCAP plus extracted features \cite{b13}; vendor
diversity there is a separate axis from process/application diversity, which we keep distinct in
Table~\ref{tab:compare}. SWaT \cite{b7} and WADI \cite{b8} remain the most widely used physical ICS
benchmarks outside the S7/Modbus family. The 2010 Stuxnet incident remains the standard motivating
reference for treating ICS network security as distinct from IT security \cite{b9}.

We did not identify an equivalent audited account of collection-induced leakage
(Section~\ref{sec:audit}) in the compared publications; our contribution is the equivalent physical-S7comm
role with process-state coupling, a symmetric multi-model evaluation-protocol comparison, and that audit.

\begin{table}[t]
\caption{Dataset comparison}
\label{tab:compare}
\centering
\small
\begin{tabular}{lccccc}
\toprule
Dataset & Phys. & S7 & PCAP & Proc.\ vars. & Sess.-disj. \\
\midrule
QUT\_S7comm \cite{b3}     & Y & Y   & Y & Y & N \\
Incinerator \cite{b4}     & Y & Y   & Y & N & N \\
Electra \cite{b5}         & Y & Y+M & Y & N & N \\
Dehlaghi et al.\ \cite{b6}& Y & N   & Y & Y & N.r. \\
Kellerer et al.\ \cite{b10}& Y& Y   & Y & Y & N.r. \\
ICS-MMOC3 \cite{b12}      & Y & Y   & Y & Y & N.r. \\
ICS-NAD \cite{b13}        & Y, 3 vend. & Part. & Y & N & N.r. \\
\textbf{This work}        & Y & Y   & Y & Y & Y (D6, expl.) \\
\bottomrule
\end{tabular}

\vspace{2pt}
\raggedright
\footnotesize N.r.\ = not reported in the source publication; M = Modbus TCP; Sess.-disj.\ =
session-disjoint evaluation split; expl.\ = exploratory rather than confirmatory
(Section~\ref{sec:audit}).
\end{table}

\section{Testbed and Data Collection}

\subsection{Architecture}

The testbed deploys a physical Siemens S7-1500 PLC (.211) executing real control logic, connected to a
genuine engineering/HMI station (TIA Portal / WinCC Runtime V18, .31) over an isolated, segmented network.
An attacker host (.32) runs Python, Scapy, and python-snap7. All hosts connect through an Aruba switch
(.254) configured for ERSPAN GRE mirroring to a dedicated capture/logging host, which independently
records network capture (PCAP), the PLC tag/process log, and the attack-event timeline used as audit
ground truth; the three streams are merged by timestamp (Fig.~\ref{fig:testbed}). A physical PLC, rather
than PLCSIM, improves the realism of cyclic scan timing and protocol-level responses under normal and
adversarial conditions. Two control applications are instrumented and released---an urban traffic-light
controller and a conveyor-belt sequential-logic process. The current release does not support an
auditable per-application window breakdown because application identifiers were not preserved at window
level (Section~\ref{sec:limitations}). A third (irrigation/water-pump) is implemented on the same PLC but
not yet captured, and is not part of this benchmark.

\textbf{Threat model.} We assume an attacker with existing network access to the Field OT segment reaching
the PLC's S7 service on TCP/102; we do not simulate initial compromise. Out of scope: physical sabotage,
bypassing a Safety PLC, firmware persistence, functional-safety evaluation. All experiments run on
authorized equipment with dedicated test memory areas and a restore step after each write episode.

\textbf{Observability and host identity.} The WinCC-to-PLC channel can use S7CommPlus, which may not be
decodable from the mirrored capture; for protected sessions only transport metadata is available. All
attack traffic originates from a single attacker host (.32), distinct from the HMI host (.31); although
host addresses are excluded from the feature matrix, directional and connection-pattern statistics could
still let a model partly learn attacker-role rather than attack behavior in general. We did not run a
cross-attacker-host evaluation to bound this (Section~\ref{sec:limitations}).

\textbf{Clock synchronization.} Streams are merged by nearest recorded timestamp; inter-host clock
synchronization (NTP/PTP) has not been confirmed, and the process-logger polling interval and the exact
merge tolerance have not been measured. Alignment error at the sub-window scale therefore cannot be
bounded from the current release. See Section~\ref{sec:limitations} for the full caveat.

\begin{figure}[t]
\centering
\fbox{\parbox{0.9\linewidth}{\centering\small
\textbf{Testbed architecture (Fig. 1 placeholder)}\\[4pt]
S7-1500 PLC (.211) $\rightarrow$ Aruba switch (.254), ERSPAN GRE mirror $\rightarrow$
Capture/logging host (TShark, Npcap)\\
Engineering/HMI station (.31, TIA Portal/WinCC) $\rightarrow$ switch\\
Attacker host (.32, Scapy, python-snap7) $\rightarrow$ switch; also emits attack-event
timeline (ground truth)\\
PLC $\rightarrow$ PLC tag/process log (register poll) $\rightarrow$ capture host\\[4pt]
All three streams merge by timestamp into 2-second windows.\\[2pt]
\textit{Replace this box with the original vector figure; enlarge label font relative to the
v3/v28 draft, which reviewers flagged as too small to read at print size.}
}}
\caption{Testbed architecture. PLC, HMI, and attacker traffic are ERSPAN-mirrored to a capture
host that independently records three streams, merged by timestamp into 2-second windows.}
\label{fig:testbed}
\end{figure}

\subsection{Definitions}

A window is a fixed 2-second interval, stride 2\,s, no overlap, labeled ATTACK if any attack interval
overlaps it. An episode is one warm-up$\rightarrow$attack$\rightarrow$cooldown$\rightarrow$restore
sequence, the grouping unit for cross-validation. A session is all data from one collection day. The
effective grouping key is \texttt{session\_id|episode\_id}; the codebase's group configuration also lists
\texttt{host\_id}, retained in the released metadata, but it does not affect grouping in this release,
since the testbed has a single fixed attacker host and a single HMI host (Section~III-A). For attack
windows, \texttt{episode\_id} identifies one of the 36 attack episodes. Benign windows are not one group
per day: they are chunked into 10-minute \texttt{episode\_id} segments
(\texttt{session\_id:BENIGN:<chunk>}), giving 192 benign groups and 228 groups in total for the network
and fusion views (49 for process-only, which contains only logger-covered rows). We call this
episode-grouped cross-validation rather than session-grouped: repeated episodes of a scenario
(Section~III-C) share script and parameters, so grouping by episode prevents seeing the exact same episode
in both folds, but not a different episode of the same scenario, script, and day. Table~\ref{tab:main}
(episode-grouped) and the Day-6 result (which additionally shifts day, order, and rate) should not be
conflated.

\subsection{Collection Schedule}

Table~\ref{tab:byday} reports the measured window count by day and attack class, queried directly from
the released feature files. Day 1 is a benign baseline; Days 2--5 each exercise a subset of scenarios
(reconnaissance-to-impact progression: enumeration/scan, then writes, then setpoint/spoof/stealthy logic
attacks, then volume attacks); Day 6 is the session-disjoint evaluation set, containing all nine scenarios
at reduced rate. FLOOD aggregates two scenario IDs, S7\_FLOOD (623 windows, S7comm connection exhaustion)
and SYN\_FLOOD (588 windows, TCP SYN flood), sharing one label.

Each of the nine attack scenarios runs three episodes across Days 2--5 (randomized duration and
inter-episode gap) plus one additional episode on Day 6 at 20--50\% lower rate, for four episodes per
scenario and 36 attack episodes in total. Benign traffic is organized into 192 benign segments (grouping
units used the same way as episodes for cross-validation, but without an attack execution phase), for 228
grouping units overall. The six sessions span approximately 49.8 wall-clock hours from the first to the
last recorded event; the released non-overlapping 2-second windows represent 31.1 of those hours (55{,}902
windows $\times$ 2\,s). Windows were generated only while both packet capture and the labeled session
timeline were active; overnight periods and intervals with capture stopped account for the remaining 18.7
hours, and no interval was excluded based on its attack/benign label or feature values. Rows within a
configurable transition margin around attack boundaries
(\texttt{-{}-drop-transition-seconds}, default 10\,s) are additionally dropped to avoid ambiguous
mixed-state windows, applied symmetrically regardless of label. Benign traffic covers startup,
steady-state, shutdown, mode changes, sensor transitions, HMI polling, and legitimate engineering writes,
retained deliberately as hard negatives.

\begin{table}[t]
\caption{Measured windows by day and class (queried from the release-candidate feature files)}
\label{tab:byday}
\centering
\scriptsize
\setlength{\tabcolsep}{3pt}
\begin{tabular}{lrrrrrrrrr}
\toprule
Day & BEN. & ENUM & FLD & FUZZ & RWR & SCAN & SETP & SPF & STL \\
\midrule
1 (base)  & 7{,}348  & 0   & 0   & 0   & 0   & 0   & 0   & 0   & 0 \\
2 (recon) & 8{,}761  & 935 & 0   & 0   & 0   & 438 & 0   & 0   & 0 \\
3 (write) & 11{,}753 & 0   & 0   & 0   & 866 & 0   & 0   & 0   & 0 \\
4 (logic) & 5{,}208  & 0   & 0   & 0   & 0   & 0   & 829 & 943 & 442 \\
5 (volume)& 9{,}787  & 0   & 894 & 410 & 0   & 0   & 0   & 0   & 0 \\
6 (eval.) & 4{,}603  & 470 & 317 & 122 & 365 & 462 & 253 & 299 & 397 \\
\midrule
\textbf{Total} & \textbf{47{,}460} & \textbf{1{,}405} & \textbf{1{,}211} & \textbf{532} &
\textbf{1{,}231} & \textbf{900} & \textbf{1{,}082} & \textbf{1{,}242} & \textbf{839} \\
\bottomrule
\end{tabular}

\vspace{2pt}
\raggedright\footnotesize Grand total: 55{,}902 windows (47{,}460 benign, 8{,}442 attack).
\end{table}

\subsection{Feature Layers and Leakage Controls}

Features are organized into four cumulative layers: \textbf{L0}---network volume/protocol statistics (100
features: packet/byte rates, TCP flags, length and payload statistics, port statistics);
\textbf{L1}---ICS protocol presence (+22: S7comm packet counts, COTP/TPKT/DCP indicators, directional PLC
traffic shares); \textbf{L2}---S7 operation semantics (+15: read/write counts, CPU-control ops, merker
writes, unique DB counts, write-ratio features; addresses mapped to process-meaningful tags, e.g.\
M5.0$\rightarrow$START); \textbf{L3}---process state (+66: \texttt{proc\_\_*} register-polling features).
The process-only view is restricted to the 10{,}771 rows with a real, logger-covered process observation
(\texttt{proc\_data\_valid}=1); the fusion view instead retains every network window and fills
\texttt{proc\_\_*} values for windows without a real observation via the per-day backward-fill in
Section~\ref{sec:audit}, so the backward-fill risk described there applies to fusion, not to process-only.

Session/day identifiers, filenames, host identity, timestamps, and rule-derived flags are excluded from
the feature matrix. \texttt{scenario\_id}/\texttt{episode\_id} are used only for grouped-CV splitting and
ground-truth labeling, never as model inputs; the attack-event log is used only for ground truth and
audit. Our merge pipeline also emits \texttt{proc\_data\_valid}, marking whether the process logger was
running for a window. Because the logger was started manually, this mask correlates with attack class as
a collection-timing artifact rather than a security signal; we exclude it explicitly as an
indirect-leakage risk.

\section{Experimental Setup}

Binary detection (benign vs.\ attack) is the primary task; multiclass attribution is secondary.
Process-state features are contemporaneous within each window, so the process and fusion views address
detection of an ongoing or already-realized attack, not prediction before physical impact. We evaluate
Random Forest, Logistic Regression, XGBoost, and CatBoost on three feature views (network-only,
process-only, fusion) under three protocols: (1) episode-grouped cross-validation
(\texttt{StratifiedGroupKFold}, 5 folds, grouping as defined in Section~III-B) on Days 1--5; (2) the Day-5
temporal holdout, training on Days 1--4 and testing on Day 5; (3) the Day-6 evaluation session
(randomized order, reduced rate), never used for threshold tuning or feature selection. Results are means
over five seeds (42--46). We report Macro-F1 as the primary comparison metric throughout
(Table~\ref{tab:main}), with MCC reported alongside it for the Day-5 holdout (Table~\ref{tab:day5}), the
layer-wise ablation (Table~\ref{tab:ablation}), and the process-only diagnosis
(Section~\ref{sec:audit}).

\textbf{Preprocessing.} NaN and infinite values are replaced with zero prior to fitting. Constant features
and features with pairwise Pearson $|r|>0.98$ are removed per training fold only, so no test-fold
information influences feature selection. Logistic Regression is preceded by a \texttt{StandardScaler};
tree-based models receive raw features.

\begin{table}[t]
\caption{Model hyperparameters}
\label{tab:hparams}
\centering
\scriptsize
\setlength{\tabcolsep}{3pt}
\begin{tabular}{lcccc}
\toprule
Param. & RF & XGBoost & CatBoost & LogReg \\
\midrule
n\_estim./iter.   & 300 & 300 & 300 & --- \\
max\_depth        & None & 6 & 6 & --- \\
learn.\ rate      & --- & 0.1 & 0.1 & --- \\
subsample         & --- & 0.8 & --- & --- \\
colsample         & --- & 0.8 & --- & --- \\
class weight      & balanced\_subsample & --- & auto (Balanced) & balanced \\
min\_samples\_leaf& 2   & --- & --- & --- \\
solver/penalty    & --- & --- & --- & lbfgs / l2 (default) \\
C                 & --- & --- & --- & 1.0 (default) \\
max\_iter         & --- & --- & --- & 2000 \\
eval\_metric      & --- & logloss/mlogloss & --- & --- \\
\bottomrule
\end{tabular}

\vspace{2pt}
\raggedright\footnotesize All models: 5 seeds (42--46), \texttt{n\_jobs}/\texttt{thread\_count}=$-1$.
LR solver/penalty/C are scikit-learn defaults (not explicitly overridden). XGBoost's objective is set
automatically by the fitted classifier (\texttt{binary:logistic} for binary, \texttt{multi:softprob} for
multiclass), with no \texttt{scale\_pos\_weight}, so it is not class-balanced in the same way as the other
three models. A sensitivity check with \texttt{scale\_pos\_weight} set to the negative/positive ratio on
each training fold is planned; results will be added as two additional rows in Table~\ref{tab:main} or
as a footnote once available. CatBoost infers its loss from the post-encoding class count and applies
\texttt{auto\_class\_weights=Balanced} in both tasks.
\end{table}

The binary decision threshold reported in Tables~\ref{tab:main}--\ref{tab:perclass} is selected using
grouped out-of-fold predictions (\fbetaoof): for each outer split, an inner grouped CV within the
outer-training partition refits constant/correlation feature pruning independently on each
inner-training fold and predicts the corresponding inner-validation fold to produce out-of-fold scores;
$F_\beta$ ($\beta=2$, favoring recall) is maximized over those scores, the model is refit on the full
outer-training partition, and the resulting threshold is applied once to the held-out partition. An
earlier in-sample variant, which selected the threshold from each model's own training-fold scores, gave
visibly different numbers (network CatBoost Day-6 Macro-F1 0.643, vs.\ 0.659$\pm$0.006 under
\fbetaoof) and a different model ranking, so we report the \fbetaoof\ results throughout unless noted
otherwise. Alarm burden is \texttt{fpr\_per\_hour} = false\_positive\_windows / benign\_hours, at window
level: 30 consecutive false-positive windows count as 30 alarms, not one episode; an episode-collapsed
variant is left to future work. Grouped-CV metrics are averaged across the five folds within each seed,
then summarized across the five seeds.

\textbf{Status of Day 6.} An early pipeline version produced an anomalous, worse-than-chance score on Day
6 (Section~\ref{sec:audit}); that anomaly triggered the data-quality audit and subsequent fixes described
there. Day-6 labels were never used to tune features or thresholds, but its results did inform a pipeline
revision, so we report Day 6 as a session-disjoint evaluation set that also served as an exploratory
diagnostic, not a confirmatory holdout. The episode-grouped result on Days 1--5 is the closest
confirmatory in-distribution result this paper offers.

Dataset composition (55{,}902 windows: 47{,}460 benign, 8{,}442 attack; per-day and per-class breakdown in
Table~\ref{tab:byday}) was cross-checked against the data card and both network-only and fusion exports;
all agree exactly.

\section{Results}

\subsection{Data Quality Audit}
\label{sec:audit}

We report this first because it governs the interpretation of everything after it, and because it is why
Day 6 is not described as an untouched holdout. An early pipeline version produced MCC $=-0.154$ (Random
Forest, process-only, Day 6)---a worse-than-chance, consistently-signed result that motivated an audit of
labels and preprocessing rather than being accepted at face value. The audit found: (1) the
\texttt{proc\_data\_valid} availability mask was computed after ffill/bfill had already overwritten
missing values, so it no longer reflected real observation coverage; (2) forward/backward fill ran across
the whole dataset in timestamp order, so Day-1 process values propagated into Day-2 attack windows
collected before the logger started; (3) the mask itself was inadvertently left in the feature set,
letting models learn \texttt{proc\_data\_valid}=0 $\rightarrow$ attack---a collection-timing artifact, not
process behavior.

The root cause is operational: the logger was started manually, 13 minutes after the last Day-2 scan
window, 36 minutes after the Day-3 write window, and only 10 seconds before the Day-5 attack period
ended---so RWRITE, FLOOD, and SETPOINT\_ATTACK have zero windows with a real process observation
(\texttt{proc\_data\_valid}=1) anywhere in their collection day. After switching
\texttt{proc\_data\_valid} assignment to precede imputation and restricting ffill/bfill to within each
day, process-only MCC moved from $-0.154$ to $+0.374$ (RF), 0.000 to $+0.376$ (LR), and $-0.007$ to
$+0.374$ (CatBoost); all subsequent results use this revised pipeline. This closes the cross-day
propagation and the mask-as-feature leak, but per-day backward-fill still assigns
RWRITE/FLOOD/SETPOINT\_ATTACK windows a process value copied from the nearest later same-day observation
(e.g.\ after logger start on Day 3), since those windows have no real observation to fall back to. We
therefore do not attempt a matched comparison restricted to \texttt{proc\_data\_valid}=1 windows as a way
to validate fusion: three attack classes have no such windows at all, and among the classes that do (SCAN
27.7\%, ENUMERATION 30.4\%, FUZZ 24.1\%, SPOOF 4.4\%), the Q0 register value is 41---identical to
BENIGN---so a proc=1-only evaluation would reduce to STEALTHY-only detection rather than a genuine
multiclass comparison (Section~\ref{sec:limitations}). Closing this requires re-collecting with the
logger running for the full duration of every session, not a different evaluation protocol on the current
data.

Under \fbetaoof, grouped-CV FPR/hour rises sharply relative to the earlier in-sample threshold (e.g.\
network CatBoost 105 vs.\ 37), while Day-6 FPR/hour stays below 3 for every network and fusion model
(1.02--2.97 for network-only, 0.08--0.39 for fusion): a less-overfit threshold fires more often on
grouped-CV benign traffic it has not memorized, but still rarely fires on Day 6, where recall is already
low (Table~\ref{tab:main}). Since the threshold optimizes $F_\beta$ rather than Macro-F1, a low Day-6
FPR/hour should not be read as good calibration; it is consistent with under-detection on the harder
split (Section~\ref{sec:limitations}). Process-only carries by far the heaviest grouped-CV burden
($\sim$720/hour for all four models), consistent with the single-register decision rule below.

\subsection{In-Distribution vs.\ Day-6 (RQ3)}

Under \fbetaoof, network-only in-distribution performance is fairly uniform (0.833--0.909 F1), and the
Day-6 gap is again model-dependent, but the ranking differs from the in-sample-threshold results: XGBoost
now shows the largest drop ($-0.255$) rather than Random Forest, with CatBoost and Random Forest close
behind ($-0.250$, $-0.239$). This reordering shows that which model looks most robust to this shift
depends on threshold calibration as well as on the model itself, so a single-model, single-threshold
report of ``domain-shift robustness'' is not reliable evidence about the shift. We therefore report the
full model$\times$protocol matrix rather than each split's best model.

All four process-only models return essentially the same Day-6 score (0.547--0.548) regardless of view or
threshold-calibration method. Feature importances show \texttt{proc\_q0\_raw\_hex\_mean} dominates every
model; the decision rule reduces to a threshold on the Q0 output register, and computing binary Macro-F1
under (Q0=40$\rightarrow$attack) by hand reproduces this value almost exactly---the process-only view
detects one register transition, not process behavior in general. Under \fbetaoof, process-only
grouped-CV Macro-F1 is markedly lower (0.389--0.402) than under the in-sample threshold (0.584--0.696): a
threshold calibrated out-of-sample on this small (10{,}771-row), heavily class-imbalanced view generalizes
far worse within grouped CV than one tuned in-sample, which is itself informative about how fragile the
process-only signal is outside the single Q0 rule that survives on Day 6.

Day 5 (11{,}091 windows, 88.2\% benign, only FLOOD and FUZZ) yields near-ceiling network and fusion scores
(RF, XGBoost $\geq$0.998 Macro-F1), exceeding every Day-6 result. A score this close to perfect warrants
scrutiny: an independent check confirmed no metadata columns remain in the feature matrix, zero
train/test group overlap, and no exact-duplicate rows between Days 1--4 and Day 5; the near-ceiling scores
instead trace to strong volume separation (e.g.\ \texttt{tcp\_syn\_count} tops out at 12 for benign
Day-5 windows vs.\ 18+ for attack windows), consistent with FLOOD/SYN-flood signatures being trivially
separable by volume. This is not stronger temporal generalization: Day 6 mixes all nine scenarios
including low-footprint attacks, and attack composition, not the schedule shift, is what makes Day 6
hard. Process-only collapses to chance (Macro-F1$\approx$0.499) on Day 5, since FLOOD/FUZZ windows have
\texttt{proc\_data\_valid}=0 (Section~\ref{sec:audit})---no signal to exploit at all, unlike the partial
STEALTHY-driven signal on Day 6.

\subsection{Layer-Wise Ablation (RQ2)}

For CatBoost on Day 6 under \fbetaoof, ICS-presence (L1) gives a small gain over volume-only
(0.648$\rightarrow$0.654 Macro-F1); S7-operation semantics (L2) does not extend this (0.647). Adding
process state (L3) gives the highest score here (0.663 Macro-F1, 0.483 MCC) alongside the lowest FPR/hour
(0.08). We do not read this as a general fusion endorsement: the ablation trains on Days 1--4 rather than
Table~\ref{tab:main}'s Days 1--5, and the process contribution concentrates in the STEALTHY-specific Q0
signal (Section~\ref{sec:audit}) rather than a broad improvement, consistent with the model-dependent
fusion effect below.

\subsection{Fusion and Multiclass Attribution}

Fusion appears to beat network-only at multiclass (0.235 vs.\ 0.206). A perfect F1=1.000 on any class in
a Day-6 split warrants suspicion, so we investigated. Of 8{,}442 attack windows, 1{,}698 (20.1\%) have
\texttt{proc\_data\_valid}=1: STEALTHY 839 (100\%), ENUMERATION 427 (30.4\%), SCAN 249 (27.7\%), FUZZ 128
(24.1\%), SPOOF 55 (4.4\%), and none for RWRITE, FLOOD, SETPOINT\_ATTACK---so availability alone does not
explain STEALTHY's perfect score.

Checked specifically on Day 6, \texttt{proc\_q0\_raw\_hex\_mean} takes the single value 40.0 across all
397 STEALTHY windows, while ENUMERATION, SCAN, and FUZZ windows with \texttt{proc\_data\_valid}=1 all
read 41.0---indistinguishable from BENIGN; only STEALTHY's STOP-bit write changes Q0. So of the 20.1\% of
attack windows with a real process observation, only the 9.9\% that are STEALTHY carry a signal that
discriminates anything at all. Removing STEALTHY reverses the ordering: fusion Macro-F1 0.139 vs.\
network-only 0.147.

Naive concatenation of network and process views does not yield a consistent binary gain on Day 6 under
\fbetaoof\ (Table~\ref{tab:main}): lower than network-only for CatBoost (0.626 vs.\ 0.659) and Random
Forest (0.634 vs.\ 0.662), higher for XGBoost (0.673 vs.\ 0.649), close for Logistic Regression. This
model-dependence, not a uniform ``fusion does not help,'' is what the corrected calibration supports.
With 79.9\% of attack windows lacking real process observations and the multiclass advantage traced
entirely to the STEALTHY artifact, this experiment cannot answer RQ2 conclusively either way; any fused
score should be audited per model and per class before being believed.

\begin{table*}[t]
\caption{Binary detection under grouped out-of-fold threshold calibration (\texttt{fbeta\_oof}), mean
$\pm$ std over five seeds\textsuperscript{a}}
\label{tab:main}
\centering
\small
\begin{tabular}{llcccccccc}
\toprule
View & Model & Grp.\ CV Macro-F1 & Day-6 Macro-F1 & $\Delta$F1 & Grp.\ CV MCC & Day-6 MCC &
\multicolumn{2}{c}{Day-6 attack\textsuperscript{b}} \\
\cmidrule(lr){8-9}
 & & & & & & & P & R \\
\midrule
network & CatBoost      & 0.909$\pm$0.002 & 0.659$\pm$0.006 & $-$0.250 & 0.831$\pm$0.003 & 0.477$\pm$0.007 & 0.996 & 0.309 \\
network & Log.\ Reg.    & 0.833$\pm$0.011 & 0.595$\pm$0.000 & $-$0.239 & 0.702$\pm$0.013 & 0.396$\pm$0.000 & 0.997 & 0.229 \\
network & XGBoost       & 0.903$\pm$0.001 & 0.649$\pm$0.006 & $-$0.255 & 0.821$\pm$0.002 & 0.463$\pm$0.007 & 0.994 & 0.304 \\
network & Rand.\ Forest & 0.900$\pm$0.003 & 0.662$\pm$0.005 & $-$0.239 & 0.817$\pm$0.005 & 0.479$\pm$0.005 & 0.996 & 0.313 \\
fusion\textsuperscript{c} & CatBoost      & 0.922$\pm$0.003 & 0.626$\pm$0.005 & $-$0.297 & 0.854$\pm$0.005 & 0.436$\pm$0.007 & 1.000 & 0.271$\pm$0.007 \\
fusion\textsuperscript{c} & Log.\ Reg.    & 0.910$\pm$0.002 & 0.621$\pm$0.018 & $-$0.289 & 0.830$\pm$0.003 & 0.430$\pm$0.022 & 0.999 & 0.265$\pm$0.022 \\
fusion\textsuperscript{c} & XGBoost       & 0.919$\pm$0.003 & 0.673$\pm$0.010 & $-$0.246 & 0.847$\pm$0.006 & 0.495$\pm$0.012 & 0.999 & 0.341$\pm$0.013 \\
fusion\textsuperscript{c} & Rand.\ Forest & 0.918$\pm$0.003 & 0.634$\pm$0.022 & $-$0.284 & 0.847$\pm$0.005 & 0.446$\pm$0.028 & 0.999 & 0.283$\pm$0.028 \\
process\textsuperscript{c} & CatBoost      & 0.389$\pm$0.027 & 0.547$\pm$0.000 & $+$0.158 & 0.116$\pm$0.043 & 0.370$\pm$0.003 & --- & --- \\
process\textsuperscript{c} & Log.\ Reg.    & 0.402$\pm$0.000 & 0.548$\pm$0.000 & $+$0.146 & 0.141$\pm$0.001 & 0.376$\pm$0.000 & --- & --- \\
process\textsuperscript{c} & Rand.\ Forest & 0.393$\pm$0.018 & 0.547$\pm$0.000 & $+$0.155 & 0.121$\pm$0.025 & 0.372$\pm$0.001 & --- & --- \\
process\textsuperscript{c} & XGBoost       & 0.401$\pm$0.001 & 0.548$\pm$0.000 & $+$0.146 & 0.135$\pm$0.004 & 0.374$\pm$0.000 & --- & --- \\
\bottomrule
\end{tabular}

\vspace{2pt}
\raggedright\footnotesize
\textsuperscript{a}Threshold selected via grouped out-of-fold predictions within each outer-training
partition (Section IV). \textsuperscript{b}Attack-class precision/recall at the \texttt{fbeta\_oof} threshold on Day~6,
computed from per-seed binary confusion matrices. Network P/R values are from a single-seed run
(seed~42); fusion P/R values are means over five seeds (std shown when $>$0.001). Process-only
P/R pending re-run of the full pipeline (threshold exists; confusion matrix with fbeta\_oof not
yet saved for this view). \textsuperscript{c}Fusion and process-only views are exploratory: process-logger
coverage is incomplete, and same-day backward-filling in the fusion view (Sections III-D
and~\ref{sec:audit}) fabricates process values for windows with no real observation; process-only excludes
those windows entirely but is correspondingly small and class-imbalanced. Process-only Grouped-CV Macro-F1
is markedly lower here than under the earlier in-sample threshold, reflecting a stricter calibration on
that small training signal.
\end{table*}

\begin{table}[t]
\caption{FPR/hour under \texttt{fbeta\_oof} (window-level, mean $\pm$ std over five seeds)}
\label{tab:fpr}
\centering
\scriptsize
\begin{tabular}{llcc}
\toprule
View & Model & Grouped CV & Day 6 \\
\midrule
network & CatBoost      & 105.4$\pm$3.0  & 1.64$\pm$0.70 \\
network & Log.\ Reg.    & 169.5$\pm$20.8 & 1.02$\pm$0.21 \\
network & XGBoost       & 114.8$\pm$2.5  & 1.88$\pm$0.75 \\
network & Rand.\ Forest & 120.1$\pm$4.6  & 2.97$\pm$1.72 \\
fusion  & CatBoost      & 86.2$\pm$4.2   & 0.08$\pm$0.18 \\
fusion  & Log.\ Reg.    & 98.1$\pm$2.7   & 0.31$\pm$0.18 \\
fusion  & XGBoost       & 91.3$\pm$4.0   & 0.39$\pm$0.00 \\
fusion  & Rand.\ Forest & 94.3$\pm$4.5   & 0.39$\pm$0.00 \\
process & CatBoost      & 720.8$\pm$1.1  & 9.11$\pm$6.07 \\
process & Log.\ Reg.    & 719.7$\pm$0.1  & 0.00$\pm$0.00 \\
process & Rand.\ Forest & 721.5$\pm$0.9  & 5.46$\pm$1.36 \\
process & XGBoost       & 721.2$\pm$1.4  & 3.04$\pm$0.00 \\
\bottomrule
\end{tabular}
\end{table}

\begin{table}[t]
\caption{Day-5 temporal holdout under \texttt{fbeta\_oof} (train D1--4, test D5)\textsuperscript{a}}
\label{tab:day5}
\centering
\scriptsize
\begin{tabular}{llccc}
\toprule
View & Model & Macro-F1 & MCC & FPR/hour \\
\midrule
network & CatBoost      & 0.993$\pm$0.011 & 0.986$\pm$0.022 & 0.44$\pm$0.21 \\
network & Log.\ Reg.    & 0.882$\pm$0.000 & 0.786$\pm$0.001 & 1.14$\pm$0.20 \\
network & Rand.\ Forest & 0.999$\pm$0.001 & 0.998$\pm$0.001 & 0.99$\pm$0.50 \\
network & XGBoost       & 0.998$\pm$0.001 & 0.996$\pm$0.001 & 1.55$\pm$0.44 \\
fusion  & CatBoost      & 0.930$\pm$0.064 & 0.874$\pm$0.115 & 0.11$\pm$0.10 \\
fusion  & Log.\ Reg.    & 0.884$\pm$0.000 & 0.790$\pm$0.000 & 0.00$\pm$0.00 \\
fusion  & Rand.\ Forest & 1.000$\pm$0.000 & 1.000$\pm$0.000 & 0.18$\pm$0.00 \\
fusion  & XGBoost       & 1.000$\pm$0.000 & 1.000$\pm$0.000 & 0.18$\pm$0.00 \\
process & CatBoost      & 0.499$\pm$0.001 & $-$0.001$\pm$0.002 & 2.41$\pm$5.39 \\
process & Log.\ Reg.    & 0.499$\pm$0.000 & 0.000$\pm$0.000 & 0.00$\pm$0.00 \\
process & Rand.\ Forest & 0.499$\pm$0.000 & 0.000$\pm$0.000 & 0.00$\pm$0.00 \\
process & XGBoost       & 0.499$\pm$0.000 & 0.000$\pm$0.000 & 0.00$\pm$0.00 \\
\bottomrule
\end{tabular}

\vspace{2pt}
\raggedright\footnotesize\textsuperscript{a}Train Days 1--4, test Day 5, threshold calibrated with the
same grouped out-of-fold procedure as Table~\ref{tab:main}; not a slice of that table, since the
train/test partition differs (Section III-C).
\end{table}

\begin{table}[t]
\caption{Cumulative feature layers, CatBoost, Day 6, \texttt{fbeta\_oof}; mean $\pm$ std over five
seeds\textsuperscript{a}}
\label{tab:ablation}
\centering
\small
\begin{tabular}{clccc}
\toprule
Cfg & Layers & Macro-F1 & MCC & FPR/hour \\
\midrule
A & L0             & 0.648$\pm$0.008 & 0.462$\pm$0.010 & 1.17$\pm$0.48 \\
B & L0+L1          & 0.654$\pm$0.010 & 0.470$\pm$0.013 & 1.25$\pm$0.64 \\
C & L0+L1+L2       & 0.647$\pm$0.011 & 0.461$\pm$0.012 & 1.17$\pm$1.11 \\
D & L0+L1+L2+L3    & 0.663$\pm$0.020 & 0.483$\pm$0.024 & 0.08$\pm$0.18 \\
\bottomrule
\end{tabular}

\vspace{2pt}
\raggedright\footnotesize\textsuperscript{a}Trained on Days 1--4; a feature-layer diagnostic, not a
slice of Table~\ref{tab:main} (Days 1--5). Uses the same leakage-exclusion list as the main pipeline,
including explicit removal of \texttt{proc\_data\_valid}.
\end{table}

\begin{table}[t]
\caption{Per-class P/R/F1, Day 6, CatBoost}
\label{tab:perclass}
\centering
\scriptsize
\setlength{\tabcolsep}{2.5pt}
\begin{tabular}{lcccccccc}
\toprule
Class & Net P & Net R & Net F1 & Fus.\textsuperscript{a} P & Fus.\textsuperscript{a} R &
Fus.\textsuperscript{a} F1 & Supp. \\
\midrule
BENIGN    & 0.701 & 0.998 & 0.824 & 0.687 & 1.000 & 0.814 & 4{,}603 \\
STEALTHY  & 0.641 & 0.736 & 0.683 & 1.000 & 1.000 & 1.000 & 397 \\
SPOOF     & 0.290 & 0.064 & 0.104 & 0.357 & 0.064 & 0.107 & 299 \\
SETPOINT  & 1.000 & 0.051 & 0.096 & 1.000 & 0.050 & 0.095 & 253 \\
SCAN      & 0.141 & 0.052 & 0.076 & 0.160 & 0.009 & 0.016 & 462 \\
FLOOD     & 1.000 & 0.038 & 0.072 & 1.000 & 0.037 & 0.066 & 317 \\
RWRITE    & 0.000 & 0.000 & 0.000 & 0.189 & 0.009 & 0.018 & 365 \\
ENUM.     & 0.000 & 0.000 & 0.000 & 0.000 & 0.000 & 0.000 & 470 \\
FUZZ      & 0.000 & 0.000 & 0.000 & 0.000 & 0.000 & 0.000 & 122 \\
\midrule
Macro     & 0.419 & 0.215 & 0.206 & 0.488 & 0.241 & 0.235 & 7{,}288 \\
\bottomrule
\end{tabular}

\vspace{2pt}
\raggedright\footnotesize\textsuperscript{a}Fusion is exploratory (Section~\ref{sec:audit}); STEALTHY's
F1=1.000 reflects a collection artifact, not detection improvement. Multiclass predictions use standard
argmax over predicted class probabilities, not a binary decision threshold, so the \texttt{fbeta\_oof}
calibration in Section IV does not apply to this table.
\end{table}

\section{Limitations}
\label{sec:limitations}

Day 6 is not a confirmatory holdout---its anomalous process-only score is what triggered the audit in
Section~\ref{sec:audit}, so it has already informed a pipeline revision. A future, genuinely frozen
evaluation session (e.g.\ a sealed Day 7) would allow a confirmatory report. ``Shift'' here means
within-testbed schedule shift: Days 1--6 share PLC, attacker host, topology, and tooling; Day 6 varies
order, rate, and gaps, not infrastructure---none of our splits are cross-testbed.

Streams are merged by nearest recorded timestamp rather than a hardware-synchronized clock; we have not
confirmed inter-host clock synchronization (NTP/PTP) or measured the clock offset between the attacker,
PLC, and capture hosts, nor documented the process-logger polling interval or the exact merge tolerance
used to assign a polled observation to a 2-second window. Alignment error at the sub-window scale
therefore cannot be bounded from the current release. This caveat applies to every claim that treats a
process observation as contemporaneous with a given 2-second window, including the Q0-register findings
in Sections~\ref{sec:audit}--D. If the polling interval and offset can be recovered from existing logs,
they should be reported as concrete bounds in a future release.

Process-logger activation and coverage were not synchronized with network capture: streams are
timestamp-aligned once recorded, but the logger's start/stop times were set manually and did not match
the capture window, so 79.9\% of attack windows have no real process observations, concentrated exactly
in the classes (RWRITE, FLOOD, SETPOINT\_ATTACK) where process impact would be most informative. A
\texttt{proc\_data\_valid}=1-only matched comparison, which would be the natural way to isolate fusion's
effect, is not merely undone but structurally uninformative on this release: those three classes have
zero such windows, and the remaining classes with real observations (SCAN, ENUMERATION, FUZZ, SPOOF) show
process feature values indistinguishable from BENIGN, so the comparison would reduce to STEALTHY-only
detection. RQ2 therefore requires re-collection with synchronized logging, not a different evaluation
protocol on the current data (Section~\ref{sec:audit}).

Attacker-host identity is confounded with the attack label: all attacks originate from one host, which
network features could partly exploit as a stand-in for attack behavior; we have not run a
cross-attacker-host evaluation. Multi-application coverage is not tested as a generalization axis, and
RQ1's application-level split is not yet quantified: both applications contribute windows across all six
days by collection design, and Table~\ref{tab:byday} quantifies distribution by day and attack class, but
we do not report a per-application (traffic-light vs.\ conveyor-belt) window/episode breakdown, because
the released window-level tables do not include an \texttt{application\_id} field that unambiguously
assigns every window or episode to a physical application; we also do not report a
leave-one-application-out evaluation. Future releases should add application-level metadata to support
this analysis.

Multiclass attribution is weak under schedule shift (Macro-F1 0.206--0.235, three classes at 0.000).
Episode-grouped CV is within-testbed and within-script: repeated episodes of a scenario share script and
parameters, so grouped-CV figures measure within-schedule generalization specifically; Day 6 and the
Day-5 holdout only partially address order and temporal shift. Two of three applications are instrumented;
the water-pump application is implemented but not yet captured. The threshold optimizes $F_\beta$
($\beta=2$), not Macro-F1: a recall-favoring threshold can understate what a Macro-F1-optimized or
fixed-0.5 threshold would show; Table~\ref{tab:perclass} gives per-class recall for Day 6 and
Table~\ref{tab:main} will also report binary attack precision/recall at the same threshold
(column pending script run; threshold already exists). The Day-5 holdout (Table~\ref{tab:day5}) and layer-wise ablation
(Table~\ref{tab:ablation}) use the same \fbetaoof\ calibration as Table~\ref{tab:main} but a different
train split (Days 1--4 rather than 1--5), so we report them as protocol-specific results rather than
direct slices of the main table. The multiclass per-class breakdown (Table~\ref{tab:perclass}) uses
argmax decisions, to which this binary threshold does not apply.

\section{Conclusion}

We presented a physical S7comm intrusion detection dataset---55{,}902 labeled windows aligning network
capture, process state, and attack timelines by timestamp across six days, organized through a
four-layer feature taxonomy---and a symmetric four-baseline benchmark under episode-grouped
cross-validation, a Day-5 temporal holdout, and a session-disjoint Day-6 evaluation.

\textbf{RQ1 (Coverage).} Table~\ref{tab:byday} shows all nine attack scenarios (eight released classes)
represented across the six-day schedule, with per-class attack-window counts ranging from 532 (FUZZ) to
1{,}405 (ENUMERATION) and every class present on Day 6; benign traffic spans startup, steady-state,
shutdown, mode-change, and legitimate-engineering-write conditions rather than idle traffic alone. The
release does not yet support a per-application breakdown of this coverage
(Section~\ref{sec:limitations}).

\textbf{RQ2 and RQ3.} The scientific content is in the gaps, not the peaks. Network-only in-distribution
Macro-F1 (0.83--0.91) is uniform enough across models that it predicts neither the Day-6 result nor its
ranking: the same in-distribution split separates models by less than 0.08 F1, while the Day-6 shift
penalty ranges from 0.239 to 0.255 depending on the learner and, more importantly, depending on how the
decision threshold is calibrated---an in-sample-threshold variant we discarded gave a visibly different
ranking (Section IV). Fusion shows no single direction of effect: worse than network-only for CatBoost
and Random Forest, better for XGBoost. A single-model, single-threshold report of either domain-shift
robustness or fusion benefit would therefore have been misleading in either direction; only the full
model$\times$threshold$\times$view matrix supports a defensible claim.

We traced that finding to its cause rather than reporting it as a general property of fusion. The process
logger was started manually and covered only part of each attack period; the resulting availability
pattern first corrupted our own pipeline through a mask correlated with attack class, then, after that
leak was closed, produced a perfect but meaningless F1 on the one class whose process signature survived.
Both are data-collection artifacts, and Day 6---the split whose anomalous score first surfaced
them---is reported as an exploratory, diagnostic-informed evaluation rather than a clean holdout.

\textbf{Future work.} A \texttt{proc\_data\_valid}=1-only matched comparison, our initial plan for
validating fusion, is shown in Section~\ref{sec:audit} to be uninformative here, since three classes have
no real-observation windows and the rest are indistinguishable from BENIGN outside STEALTHY. The priority
is re-collecting with the logger running for the full session and with measured inter-host clock offsets,
paired with a frozen, previously unseen evaluation session. Beyond that: leave-one-application-out and
cross-attacker-host evaluations; the water-pump application; OPC-UA/Web-SCADA channels; an episode-level
alarm-burden metric; sequence-model baselines.

\textbf{Availability.} Extracted feature views (network, process, fusion), split definitions
(\texttt{episode\_id}/\texttt{session\_id} assignments), the preprocessing code
(\texttt{merge\_dataset.py}, \texttt{train\_ml.py}), and a data card documenting window size, stride,
label precedence, and clock-offset policy are available to reviewers at submission via a linked, directly
accessible repository (no manual approval step), identified by a fixed code commit. Full PCAP/PCAPNG
release follows acceptance due to \textbf{[FILE SIZE / INSTITUTIONAL POLICY / PII --- fill in reason]};
a sanitized full-session PCAP (not only a sample) is included in the reviewer package.

\begin{thebibliography}{13}
\bibitem{b1} A. Termanini, H. Bourdoucen, D. Al-Abri, and A. Al Maashri, ``Intrusion Detection Datasets
for IIoT and ICS: A Taxonomic Review with a Decision-Aid Scoring Rubric,'' \emph{Sensors}, vol. 26, no.
13, Art. no. 4099, 2026.
\bibitem{b2} A. Kleinmann and A. Wool, ``Accurate Modeling of the Siemens S7 SCADA Protocol for Intrusion
Detection and Digital Forensics,'' \emph{J. Digital Forensics, Security and Law}, vol. 9, no. 2, pp.
37--50, 2014.
\bibitem{b3} N. R. Rodofile, T. Schmidt, S. T. Sherry, C. Djamaludin, K. Radke, and E. Foo, ``Process
Control Cyber-Attacks and Labelled Datasets on S7Comm Critical Infrastructure,'' in \emph{Information
Security and Privacy (ACISP 2017)}, LNCS vol. 10343, Springer, 2017, pp. 452--459.
\bibitem{b4} B. Al-Duwairi, A. Shatnawi, A. Al-Hammouri, and M. Ababneh, ``Dataset of SCADA traffic
captures from a medical waste incinerator with injected cyberattacks,'' \emph{Data in Brief}, vol. 63,
Art. no. 112294, 2025.
\bibitem{b5} A. L. Perales G\'omez, L. Fern\'andez Maim\'o, A. Huertas Celdr\'an, F. J. Garc\'{i}a
Clemente, C. Cadenas Sarmiento, C. J. Del Canto Masa, and R. M\'endez Nistal, ``On the Generation of
Anomaly Detection Datasets in Industrial Control Systems,'' \emph{IEEE Access}, vol. 7, pp.
177460--177473, 2019.
\bibitem{b6} A. Dehlaghi-Ghadim, M. Helali Moghadam, A. Balador, and H. Hansson, ``Anomaly Detection
Dataset for Industrial Control Systems,'' \emph{IEEE Access}, vol. 11, pp. 107982--107996, 2023.
\bibitem{b7} A. P. Mathur and N. O. Tippenhauer, ``SWaT: A Water Treatment Testbed for Research and
Training on ICS Security,'' in \emph{CySWater}, 2016, pp. 31--36.
\bibitem{b8} C. M. Ahmed, V. R. Palleti, and A. P. Mathur, ``WADI: A Water Distribution Testbed for
Research in the Design of Secure Cyber Physical Systems,'' in \emph{CySWATER'17}, ACM, 2017, pp. 25--28.
\bibitem{b9} T. M. Chen, ``Stuxnet, the real start of cyber warfare? [Editor's Note],'' \emph{IEEE
Network}, vol. 24, no. 2, pp. 2--3, 2010.
\bibitem{b10} N. Kellerer, G. Sanchez Collado, H. Alberto, V. Hagenmeyer, and G. Elbez, ``S7 Data
Modification Attacks using an Industrial Control System Testbed,'' Zenodo, May 2025. doi:
10.5281/zenodo.15373938.
\bibitem{b11} N. Kellerer, G. Sanchez, H. Alberto, V. Hagenmeyer, and G. Elbez, ``Attacks on the Siemens
S7 Protocol Using an Industrial Control System Testbed,'' in \emph{Proc. 16th ACM Int. Conf. Future and
Sustainable Energy Systems (e-Energy '25)}, 2025, pp. 770--779. doi: 10.1145/3679240.3734645.
\bibitem{b12} L. Shafir, A. Elyashar, J. Cohen, M. Zeitoun, M. Holipsky, B. Yaakov, E. Aliev, R. Puzis, Y.
Harel, and A. Wool, ``ICS-MMOC3: Exploring Operational Faults and Cyber Attacks with a Realistic
Multi-Modal Power Generation ICS Dataset,'' in \emph{Proc. 12th ACM Cyber-Physical System Security
Workshop (CPSS 2026)}, 2026, pp. 63--76. doi: 10.1145/3775042.3807883.
\bibitem{b13} X. Zhou, Z. Cheng, C. Wang, S. Wang, C. Tao, Z. Zhou, et al., ``A dataset collected in
real-world industrial control systems for network attack detection,'' \emph{Scientific Data}, vol. 13,
Art. no. 399, 2026. doi: 10.1038/s41597-026-06738-x.
\end{thebibliography}

\end{document}
